% Reserve more space for formal result blocks and less for worked material that may flow naturally.
\newcommand{\formalblockneedspace}{10\baselineskip}
\newcommand{\supportingblockneedspace}{7\baselineskip}
\newcommand{\workedblockneedspace}{4\baselineskip}
\newcommand{\proofblockneedspace}{3\baselineskip}
\newcommand{\reserveblockspace}[2]{%
  \AddToHook{env/#1/before}{\Needspace{#2}}%
}

% ---------------------------------------------------------------------------
% Shared counter group for formal statements.
% Per CONTENT_README.md (Numbering Policy), the environments
%   definition, theorem, proposition, lemma, corollary
% share a single chapter-scoped counter and increment in order of appearance.
%
% We implement this with aliascnt rather than amsthm's [theorem] optional
% argument. The [theorem] shortcut works for the printed number, but it
% also makes amsthm tag every label with counter name "theorem", which in
% turn makes cleveref print every formal-statement reference as "Theorem"
% regardless of the environment actually used. aliascnt gives each
% environment its own counter name that steps the shared master counter,
% preserving distinct cleveref types.
% ---------------------------------------------------------------------------

\theoremstyle{plain}
\newtheorem{theorem}{Theorem}[chapter]

\newaliascnt{proposition}{theorem}
\newtheorem{proposition}[proposition]{Proposition}
\aliascntresetthe{proposition}

\newaliascnt{lemma}{theorem}
\newtheorem{lemma}[lemma]{Lemma}
\aliascntresetthe{lemma}

\newaliascnt{corollary}{theorem}
\newtheorem{corollary}[corollary]{Corollary}
\aliascntresetthe{corollary}

\theoremstyle{definition}
\newaliascnt{definition}{theorem}
\newtheorem{definition}[definition]{Definition}
\aliascntresetthe{definition}

% Worked blocks keep their own chapter-scoped counters.
\newtheorem{example}{Example}[chapter]
\newtheorem{exercise}{Exercise}[chapter]

\theoremstyle{remark}
\newtheorem{remark}{Remark}[chapter]

\reserveblockspace{theorem}{\formalblockneedspace}
\reserveblockspace{proposition}{\formalblockneedspace}
\reserveblockspace{lemma}{\formalblockneedspace}
\reserveblockspace{corollary}{\formalblockneedspace}
\reserveblockspace{definition}{\formalblockneedspace}
\reserveblockspace{example}{\workedblockneedspace}
\reserveblockspace{exercise}{\workedblockneedspace}
\reserveblockspace{remark}{\supportingblockneedspace}

% ---------------------------------------------------------------------------
% Solution environment.
% Visually distinct from proof, per CONTENT_README.md (Example, Solution, and
% Proof Policy, rule 11): bold "Solution." label, upright body text, and a
% filled QED box at the end (proof uses italic "Proof." with an open box).
% Optional argument overrides the label, e.g. \begin{solution}[Alternative solution].
% ---------------------------------------------------------------------------
\newenvironment{solution}[1][Solution]{%
  \par\needspace{\workedblockneedspace}%
  \pushQED{\qed}%
  \renewcommand{\qedsymbol}{\ensuremath{\blacksquare}}%
  \medskip\noindent\textbf{#1.}\enspace\ignorespaces
}{%
  \popQED\par\medskip
}

\reserveblockspace{proof}{\proofblockneedspace}

% ---------------------------------------------------------------------------
% cleveref name bindings.
% With aliascnt, each formal-statement environment has its own counter name
% that cleveref sees, so these \crefname declarations actually take effect.
% ---------------------------------------------------------------------------
\crefname{theorem}{Theorem}{Theorems}
\Crefname{theorem}{Theorem}{Theorems}
\crefname{proposition}{Proposition}{Propositions}
\Crefname{proposition}{Proposition}{Propositions}
\crefname{lemma}{Lemma}{Lemmas}
\Crefname{lemma}{Lemma}{Lemmas}
\crefname{corollary}{Corollary}{Corollaries}
\Crefname{corollary}{Corollary}{Corollaries}
\crefname{definition}{Definition}{Definitions}
\Crefname{definition}{Definition}{Definitions}
\crefname{example}{Example}{Examples}
\Crefname{example}{Example}{Examples}
\crefname{exercise}{Exercise}{Exercises}
\Crefname{exercise}{Exercise}{Exercises}
\crefname{remark}{Remark}{Remarks}
\Crefname{remark}{Remark}{Remarks}
